
Une corde de guitare a une masse de $\v{Masse}\text{ g}$ et une longueur de $\v{Lcorde}\text{ cm}$. Une flûte, considérée comme un tuyau ouvert aux deux bouts, a une longueur de $\v{Lflute}\text{ cm}$. Lorsque la corde de guitare vibre dans son deuxième mode, en même temps que la flûte vibre dans son mode fondamental, les deux instruments produisent des sons de même fréquence.

\begin{enumerate}[label=\alph*)]
    \item Tracez le schéma du mode de vibration de la corde et de la flûte. (4 pts)
    \item Déterminez la tension dans la corde de guitare. (12 pts)
    \item Quelle est la fréquence de la corde lorsqu'elle vibre dans son mode fondamental ? Quelle est alors la longueur d'onde de l'onde stationnaire sur la corde ? (schéma requis) (6 pts)
\end{enumerate}